% Compact inline orientation graphic for the Experiment 3 transfer design: a 2x2
% grid crossing generative mechanism (rows) with surface domain (columns).
%
% Row cards carry the mechanism detail: a two-node vs three-node graph with the
% mediator drawn hollow (it is never named to the model) and a persistence loop,
% plus the six-period impulse response on labelled axes. Bar heights are the
% normalised responses to a unit shock under worlds.transition at the midpoint of
% each sampling band (direct: gain .57; mediated: gain .57, rho .61, phi .30,
% lambda .81), i.e. A = (1, 0, 0, 0, 0, 0) and B = (.81, .74, .52, .34, .21, .13).
%
% Geometry note: cards have fixed dimensions and every label is a single line
% with no text width, so nothing can rewrap into a neighbour.
% Callers may set \transfergridlines to match the paragraph they anchor this
% to; wrapfig cannot wrap a section heading, so the count must not overrun it.
\providecommand{\transfergridlines}{13}
\begin{wrapfigure}[\transfergridlines]{r}{0.48\textwidth}
\vspace{-0.8em}
\centering
\definecolor{miniCityFill}{HTML}{E8F2FB}
\definecolor{miniCityEdge}{HTML}{4F7EA8}
\definecolor{miniHarborEdge}{HTML}{3F8E82}
\definecolor{miniTestFill}{HTML}{E7F4EA}
\definecolor{miniTestEdge}{HTML}{6FA97F}
\definecolor{miniTestChip}{HTML}{5E9470}
% Coin discs reuse Figure 1's yellow pair exactly -- the fill and stroke of its
% NAIVE box -- rather than the saturated gold used before. Figure 1's other warm
% pair (#FFE6CC on #D79B00, its FREQUENTIST box) is the orange one and is not
% used here. The glyph inside each disc still carries its domain colour.
\definecolor{miniCoinFill}{HTML}{FFF2CC}
\definecolor{miniCoinRim}{HTML}{D6B656}
\definecolor{miniMute}{HTML}{8A93A3}
\definecolor{miniInk}{HTML}{3F4550}
\definecolor{miniBar}{HTML}{7C93AD}
% Coin City: a skyline stamped into a coin.
\newcommand{\coincity}[2]{%
  \begin{scope}[shift={(#1,#2)},scale=0.70]
    \fill[miniCoinFill] (0,0) circle (.37);
    \draw[miniCoinRim,line width=.75pt] (0,0) circle (.37);
    \fill[miniCityEdge] (-.23,-.18) rectangle (-.10,.04);
    \fill[miniCityEdge] (-.07,-.18) rectangle (.07,.16);
    \fill[miniCityEdge] (.10,-.18) rectangle (.24,.09);
    \draw[miniCityEdge,line width=.65pt] (-.27,-.19)--(.27,-.19);
  \end{scope}}
% Coin Harbor: a small vessel and waves stamped into a coin.
\newcommand{\coinharbor}[2]{%
  \begin{scope}[shift={(#1,#2)},scale=0.70]
    \fill[miniCoinFill] (0,0) circle (.37);
    \draw[miniCoinRim,line width=.75pt] (0,0) circle (.37);
    \fill[miniHarborEdge] (-.24,-.05)--(.24,-.05)--(.15,-.18)--(-.15,-.18)--cycle;
    \draw[miniHarborEdge,line width=.7pt]
      (-.01,-.05)--(-.01,.15)--(.15,.02)--(-.01,.02);
    \draw[miniHarborEdge,line width=.55pt]
      (-.26,-.24)..controls(-.18,-.19)and(-.10,-.29)..(-.02,-.24)
      ..controls(.06,-.19)and(.14,-.29)..(.26,-.24);
  \end{scope}}
% Impulse response on labelled axes: response (Delta y) against period (t).
\newcommand{\impulse}[3]{% #1 = axis corner x, #2 = baseline y, #3 = bar heights
  \draw[-{Stealth[length=2.6pt,width=2.2pt]},draw=miniInk!65,line width=.5pt]
    (#1,#2) -- (#1,#2+0.40);
  \draw[-{Stealth[length=2.6pt,width=2.2pt]},draw=miniInk!65,line width=.5pt]
    (#1,#2) -- (#1+0.72,#2);
  \node[anchor=south,font=\sffamily\scriptsize,text=miniInk,inner sep=1pt]
    at (#1,#2+0.40) {$\Delta y$};
  \node[anchor=west,font=\sffamily\scriptsize,text=miniInk,inner sep=1.5pt]
    at (#1+0.72,#2) {$t$};
  \foreach \i/\h in {#3}{%
    \fill[miniBar] (#1+0.055+\i*0.105,#2) rectangle (#1+0.055+\i*0.105+0.070,#2+\h);}}
\resizebox{\linewidth}{!}{%
\begin{tikzpicture}[
  x=1cm, y=1cm,
  cell/.style={rounded corners=3pt, minimum width=1.95cm, minimum height=1.25cm,
    line width=.6pt, inner sep=0pt},
  traincell/.style={cell, draw=miniCityEdge, fill=miniCityFill},
  heldcell/.style={cell, draw=miniTestEdge, fill=miniTestFill,
    dash pattern=on 1.9pt off 1.5pt},
  rowcard/.style={rounded corners=3pt, minimum width=2.60cm,
    minimum height=1.25cm, inner sep=0pt, draw=miniMute!30, fill=black!3,
    line width=.5pt},
  dname/.style={anchor=south, font=\sffamily\bfseries\scriptsize},
  rname/.style={anchor=center, font=\sffamily\scriptsize, text=miniInk},
  note/.style={anchor=center, font=\sffamily\scriptsize, text=miniInk},
  gnode/.style={circle, draw=miniInk!75, fill=white, line width=.5pt,
    inner sep=0pt, minimum size=0.26cm, font=\sffamily\tiny, text=miniInk},
  hidden/.style={gnode, draw=miniMute, fill=miniMute!14,
    dash pattern=on 1.1pt off .9pt},
  garrow/.style={-{Stealth[length=2.4pt,width=2pt]}, draw=miniInk!70,
    line width=.45pt, shorten >=.4pt, shorten <=.4pt},
  chip/.style={anchor=center, rounded corners=1.6pt, draw=none, text=white,
    inner xsep=3pt, inner ysep=1.1pt, font=\sffamily\bfseries\tiny},
]
% ---- domain marks and names ----------------------------------------------
\coincity{0.975}{2.56}
\coinharbor{3.075}{2.56}
\node[dname, text=miniCityEdge]   at (0.975,1.80) {COIN CITY};
\node[dname, text=miniHarborEdge] at (3.075,1.80) {COIN HARBOR};

% ---- Structure A: direct, one period -------------------------------------
\node[rowcard] at (-1.50,1.00) {};
\node[rname] at (-1.50,1.45) {\textbf{Structure A} direct};
\node[gnode] (ax) at (-2.51,0.82) {$x$};
\node[gnode] (ay) at (-2.07,0.82) {$y$};
\draw[garrow] (ax) -- (ay);
\impulse{-1.22}{0.62}{0/0.360, 1/0.000, 2/0.000, 3/0.000, 4/0.000, 5/0.000}

% ---- Structure B: mediated, persistent -----------------------------------
\node[rowcard] at (-1.50,-0.42) {};
\node[rname] at (-1.50,0.03) {\textbf{Structure B} mediated};
\node[gnode]  (bx) at (-2.51,-0.60) {$x$};
\node[hidden] (bm) at (-2.07,-0.60) {$m$};
\node[gnode]  (by) at (-1.63,-0.60) {$y$};
\draw[garrow] (bx) -- (bm);
\draw[garrow] (bm) -- (by);
\draw[garrow] (bm) to[out=118,in=62,looseness=4.5] (bm);
\impulse{-1.22}{-0.80}{0/0.292, 1/0.265, 2/0.188, 3/0.123, 4/0.077, 5/0.048}

% ---- the four evaluation cells -------------------------------------------
\node[traincell] at (0.975,1.00) {};
\node[heldcell]  at (3.075,1.00) {};
\node[heldcell]  at (0.975,-0.42) {};
\node[heldcell]  at (3.075,-0.42) {};

\node[chip, fill=miniCityEdge] at (0.975,1.26) {TRAINED};
\node[chip, fill=miniTestChip] at (3.075,1.26) {TEST};
\node[chip, fill=miniTestChip] at (0.975,-0.16) {TEST};
\node[chip, fill=miniTestChip] at (3.075,-0.16) {TEST};

\node[note] at (0.975,0.80) {in-distribution};
\node[note] at (3.075,0.80) {domain shift};
\node[note] at (0.975,-0.62) {structure shift};
\node[note] at (3.075,-0.62) {joint shift};
\end{tikzpicture}%
}
\vspace{-0.3em}
\caption{\footnotesize\textbf{Transfer grid.} Training uses the shaded cell;
all four are evaluated. Columns vary domain; rows vary response mechanism.}
\label{fig:exp3-transfer-grid}
\vspace{-0.9em}
\end{wrapfigure}
